\chapter{Verifying a Field: The Full Campaign}
\label{ch:field}

\section{The summit statement}

Every thread so far --- specs as types, tactics, automation, $\Fp$,
primality, extraction, denotation --- was preparation for one theorem. In
the companion repositories it is called the \emph{field implementation
certificate}, and (lightly paraphrased) it says:

\begin{lstlisting}[language=Lean]
theorem fieldImplementation :
    -- p is prime, so ZMod p is genuinely a field
    Nat.Prime p
    -- and for all bounded limb arrays, every operation's
    -- commuting square closes:
  ∧ (∀ a b, Bnd a → Bnd b → AddSquare a b)
  ∧ (∀ a b, Bnd a → Bnd b → SubSquare a b)
  ∧ (∀ a b, Bnd a → Bnd b → MulSquare a b)
  ∧ (∀ a,   Bnd a →         SquareSquare a)
  ∧ (∀ a,   Bnd a →         InvertSquare a)   -- Fermat chain: a^(p-2)
  ∧ ...                                       -- reduce, negate, encode
\end{lstlisting}

One theorem, kernel-checked, quantified over \emph{every} input the
representation admits: the extracted dalek field arithmetic implements
$\Fp$. This chapter is the story of the campaign that proves it --- told
honestly, including the two places where the terrain fought back, because
the failures teach more than the victories.

\section{Order of battle}

You do not prove such a conjunction by heroism; you prove it by sequencing.
The campaign order in the real projects, and the reason for each position:

\begin{enumerate}[leftmargin=1.6em]
\item \textbf{Bounds lemmas first} --- pure \lean{omega} facts about limb
  sizes, no denotation at all. Cheap, and everything depends on them.
\item \textbf{Primality} (Chapter~\ref{ch:prime}) --- independent of the
  code entirely; it dignifies \lean{ZMod p} into a field.
\item \textbf{add, sub, negate} --- linear operations; the commuting squares
  close with \lean{omega} plus the denotation unfolds. Confidence builders.
\item \textbf{reduce} --- the $\times 19$ fold in isolation. Proving it
  alone, before mul uses it, halves the hardest proof's size.
\item \textbf{mul, square} --- the boss fight of Chapter~\ref{ch:denotation}:
  column sums, carries, folds. Squaring is mul with algebraic shortcuts ---
  a separate code path in dalek, hence a separate theorem. No shortcuts
  in the proof: the \emph{code's} shortcut is precisely what needs checking.
\item \textbf{invert} --- the 254-squaring Fermat chain, verified as a
  \lean{calc} of exponent bookkeeping on top of \code{mul}/\code{square}
  specs, ending at $a^{p-2}$; Fermat's little theorem (Mathlib's) closes
  the square.
\end{enumerate}

Notice the shape: \emph{each layer consumes only the specs of the layer
below}, never reaching into implementations. By step 6 you are doing exponent
arithmetic, blissfully ignorant of carries. That is the compositionality that
Chapter~\ref{ch:denotation}'s two-clause specs (bounds + value) were designed
to buy.

\begin{worked}{why $16p$ --- auditing a design constant to the bit}
Step 3's subtraction hides the field layer's most quotable design
decision, and you can audit it completely. The problem: compute $a - b$
limb-wise in unsigned words, where limbs of $b$ may be as large as
$2^{54} - 1$ (the invariant's edge) while limbs of $a$ may be as small as
$0$. Unsigned subtraction would truncate (Chapter~\ref{ch:lean}); the fix
is to compute $a + (kp - b)$ for a suitable multiple $k$ of the prime ---
denotationally free, since $kp \equiv 0$ --- chosen so that every limb of
$kp - b$ stays positive. How large must $k$ be? The limb constants of
$kp$ scale the telescoped representation from Chapter~\ref{ch:denotation}:
\[
kP = \big(k(2^{51}-19),\; k(2^{51}-1),\; k(2^{51}-1),\; k(2^{51}-1),\;
k(2^{51}-1)\big),
\]
whose \emph{smallest} limb is $k(2^{51} - 19)$. Positivity of every limb
of $kp - b$ demands
\[
k\,(2^{51} - 19) \;\ge\; 2^{54} - 1 .
\]
Try the candidates, because powers of two are what the code wants:
\[
k = 8:\quad 8(2^{51}-19) = 2^{54} - 152 \;<\; 2^{54} - 1
\qquad\text{✗ fails, by 151;}
\]
\[
k = 16:\quad 16(2^{51}-19) = 2^{55} - 304 \;\ge\; 2^{54} - 1
\qquad\text{✓ with } 2^{54} - 303 \text{ to spare.}
\]
So $16$ is the \emph{least} power of two that works --- $8$ misses by a
margin of $151$, invisible to any test that doesn't drive a limb of $b$
within $151$ of the very top of its envelope. Two more checks complete
the audit: the result's limbs are bounded by
$2^{54} + 2^{55} < 2^{56} < 2^{64}$ (no overflow on the additive side ✓),
and the denotation shifts by exactly $16p \equiv 0$ ✓. In the shipped
proof this box is one lemma with three \lean{omega} obligations; in the
git history of more than one real crypto library, the analogous constant
is a patched CVE. Now you know how to tell which side of that line a
codebase is on.
\end{worked}

\begin{worked}{verifying the inversion chain's bookkeeping --- all 265 steps}
Step 6 sounds heroic --- verify a hand-crafted chain of $254$ squarings
and $11$ multiplications computes $a^{p-2}$ --- until you see that the
whole verification is \emph{exponent arithmetic}, and the exponents obey
one identity you can check at every scale. The chain builds values
$a^{e}$ for a cascade of exponents $e$; squaring doubles $e$, multiplying
adds them. The dalek chain's opening moves:
\[
\begin{array}{lcl}
z_2 = a^2 & & e = 2\\
z_9 = (z_2)^{2^2} \cdot a & & e = 2\cdot 4 + 1 = 9\\
z_{11} = z_9 \cdot z_2 & & e = 9 + 2 = 11\\
z_{2^5-1} = (z_{11})^{2} \cdot z_9 & & e = 22 + 9 = 31 = 2^5 - 1 .
\end{array}
\]
From $31 = 2^5 - 1$ the chain climbs by a doubling ladder whose single
identity you should verify once symbolically:
\[
(2^k - 1)\cdot 2^k + (2^k - 1) \;=\; 2^{2k} - 1
\qquad\text{(``shift a block of $k$ ones up by $k$ and fill'').}
\]
Check it at $k=5$: $31 \cdot 32 + 31 = 992 + 31 = 1023 = 2^{10}-1$ ✓.
The ladder then assembles mixed blocks the same way:
$2^{10}-1 \to 2^{20}-1 \to 2^{40}-1$, then
$(2^{40}-1)2^{10} + (2^{10}-1) = 2^{50}-1$, then
$2^{100}-1 \to 2^{200}-1$, then
$(2^{200}-1)2^{50} + (2^{50}-1) = 2^{250}-1$. Final move: shift five and
add the stashed $z_{11}$:
\[
(2^{250}-1)\cdot 2^{5} + 11 \;=\; 2^{255} - 32 + 11 \;=\; 2^{255} - 21
\;=\; p - 2 . \qquad ✓
\]
Every line above is hand-checkable, and together they \emph{are} the
correctness of the inversion chain --- modulo only ``squaring squares and
multiplication adds exponents,'' which is the already-proven \code{mul}
and \code{square} specs feeding Fermat (Chapter~\ref{ch:modular}). Count
the cost while you are here. Squarings, stage by stage:
$1$ (to $z_2$) $+\,2$ (to $z_9$) $+\,1$ (to $z_{31}$) $+\,5+10+20$ (to
$2^{10},2^{20},2^{40}$) $+\,10$ (to $2^{50}$) $+\,50+100$ (to
$2^{100},2^{200}$) $+\,50$ (to $2^{250}$) $+\,5$ (final shift)
$= 254$. Multiplications: one per block-join --- count the ``$\cdot$''s
above --- exactly $11$. The advertised $254 + 11$ is not folklore; you
just audited it. That is the point of this box: ``verify 265 field
operations'' collapsed into ten lines of exponent algebra a patient
undergraduate audits over coffee.
\end{worked}

\section{Dispatches from the terrain}

\textbf{The wall that was really there.} Partway up, one proof style hit a
genuine limit of the tool: correctness certificates for \code{mul}-scale
goals, when handed to a general decision procedure in one monolithic call,
generate internal certificates with coefficients on the order of $2^{256}$
--- and checking them can exhaust the proof checker's memory. One such call,
during the development of the Pasta field proofs, consumed twelve gigabytes
and crashed the machine (Chapter~\ref{ch:automation} told you this story
from the tactic side). The cure was never cleverness --- it was
\emph{decomposition}: isolate each carry step as its own small lemma with a
tiny context, prove the value identity with \lean{linear_combination}
(a tactic that checks a \emph{stated} linear certificate rather than
searching for one), and let the big theorem be an assembly of small
checked parts. The same discipline, plus hard memory caps on the checker
process, became house infrastructure.

\begin{worked}{the wall, quantified --- why the monolithic proof dies}
The memory catastrophe is not mystical; you can estimate it on one page,
and the estimate teaches the cure. When a decision procedure like
\lean{omega} proves a linear goal, it does not just say ``true'' --- it
emits a \emph{certificate}: a linear combination of the hypotheses with
explicit integer coefficients that the kernel then re-checks
(Chapter~\ref{ch:prime}'s find/check split, again). Now count what a
\emph{monolithic} field-multiplication goal hands it. The goal relates
quantities at wildly different scales: raw limbs ($\sim 2^{51}$), column
sums ($\sim 2^{111}$), folded values ($\sim 2^{116}$), and the final
positional stack, whose top weight is $2^{204}$ --- so eliminating
variables across the whole system produces certificate coefficients up
to the \emph{products} of these scales: numbers of order $2^{256}$ to
$2^{512}$, i.e.\ integers of $80$--$150$ decimal digits. Per
coefficient that is only tens of bytes --- but the elimination is
quadratic-ish in the hypothesis count: with $\sim 60$ hypotheses (five
limbs each for two inputs, nine columns, nine carries, plus bounds for
everything), the certificate and the kernel's intermediate terms
multiply into \emph{millions} of big-number nodes, each participating
in further products. Twelve gigabytes is not even surprising once you
draw the multiplication table of scales --- and this estimate, run
\emph{before} the tactic, is how the wall gets predicted rather than
hit.

The cure now reads directly off the arithmetic: cut the elimination
range. Prove each carry step as its \emph{own} lemma with only its own
few hypotheses in scope (coefficients stay near the step's own scale);
where a large linear identity is genuinely needed, state it yourself
and let \lean{linear_combination} \emph{check} your stated certificate
instead of searching for one (checking is cheap; finding at scale was
the memory bomb). The companion repos' rule --- no blanket automation
in fat contexts --- is this box, written as policy. And note the shape
of the lesson: it is the \emph{same} headroom discipline as the code's
(bound your intermediates, spend your budget consciously), applied one
level up, to the proof itself.
\end{worked}

\begin{aha}
There is a deep symmetry in that fix worth savoring: the \emph{proof} was
restructured exactly the way the \emph{code} was --- into small steps with
controlled intermediate size. Delayed carries in the implementation; small
lemmas in the verification. Bounded limbs; bounded contexts. Good proofs and
good fast code turn out to obey the same engineering aesthetics. This is not
a coincidence; both are fighting combinatorial growth with structure.
\end{aha}

\textbf{Four forks, one method, real divergence.} The companion projects
verify not just upstream \code{curve25519-dalek} but three production forks
(Solana's, RISC~Zero's, Betrusted's) --- each against \emph{its own}
extraction. Worth it? The audit found the forks implement the same
constant-time conditional selection three different ways (a \code{subtle}
crate trait, a volatile-read \code{black_box}, a hand-rolled arithmetic
mask), and one fork reorders instructions in point doubling. All correct ---
\emph{provably}, now --- but the divergence is exactly the kind of thing
that silently breaks when someone ``harmonizes'' code during a rebase.
Per-fork verification is not pedantry; it is how you notice that ``the same
library'' isn't.

\begin{pitfall}
A verified fork is verified \emph{at a commit}. Change one line of
arithmetic and the certificate is stale --- that is a feature (the proof
\emph{should} break when the code changes), but it means verification is a
\emph{process wired into maintenance}, not a trophy. The companion repos
ship \code{check.sh} scripts that re-extract and re-verify from scratch;
treat those as the project's pulse, not as CI decoration.
\end{pitfall}

\section{Reading a certificate like a professional}

Suppose a stranger hands you a repository claiming ``formally verified
field arithmetic.'' Chapter~\ref{ch:honesty} gives you the full audit
protocol, but the field-layer questions you can already ask are these:

\begin{itemize}[leftmargin=1.4em]
\item \textbf{What is denoted?} Find the denotation function. Does it map
  the \emph{extracted} representation (good) or a hand-written lookalike
  (Chapter~\ref{ch:rust}'s hole)?
\item \textbf{Is the square complete?} Value equation \emph{and} bounds
  propagation \emph{and} the \lean{.ok} clause --- a spec that assumes
  success proves nothing about overflow.
\item \textbf{Is the quantifier honest?} \lean{∀ a b} with bounds
  hypotheses, or a handful of \lean{example}s on constants dressed up as
  coverage?
\item \textbf{Does anything say \lean{sorry}?} One \lean{sorry} anywhere in
  the dependency chain and the certificate is decorative. The checker will
  tell you; ask it.
\end{itemize}

\begin{tryit}
Do the audit for real: open \code{dalek-ed25519-verified}, find the
denotation, pick the \code{sub} spec, and check the three clauses against
the list above. Then run the repo's \code{check.sh} and watch the whole
pyramid rebuild. Total time: one coffee. Skill acquired: permanent.
\end{tryit}

\section*{Exercises}

\exercise{Field subtraction in dalek computes $a - b$ as
$a + (16p - b)$ limb-wise (adding a multiple of $p$ keeps limbs positive ---
recall \lean{Nat} truncation!). State the commuting square for \code{sub}
including the exact multiple-of-$p$ fact you would need as a lemma, and
explain why $16p$ rather than $p$.}

\exercise{The Fermat inversion chain computes $a^{p-2}$ via 254 squarings
and 11 multiplications. Verify the exponent bookkeeping for the first three
steps of dalek's actual chain: $a^2$, $a^{9} = (a^2)^{2\cdot 2} \cdot a$,
$a^{11} = a^9 \cdot a^2$. (The full chain is exercise-by-induction: each
step's exponent is a sum of previous ones --- addition-chain arithmetic.)}

\exercise{(Discussion) Step 5 refuses to trust the squaring shortcut and
verifies \code{square} separately from \code{mul}. A colleague argues
``square is just \code{mul a a}, reuse the theorem.'' What, concretely,
would that argument miss about the code under verification?}

\section*{Solutions and pathways}
\solutionsintro

\solhead{10.1}
\pathway Assemble from parts you own: the two-clause spec shape
(Chapter~\ref{ch:denotation}), the multiple-of-$p$ freedom (the telescope
worked example), and the positivity audit (this chapter's $16p$ worked
example). The exercise is really asking you to \emph{compose} them into
one statement.
\answer The commuting square for \code{sub}, with the lemma named:
\[
\mathrm{Bnd54}(a) \to \mathrm{Bnd54}(b) \to
\exists c,\; \code{sub}\,a\,b = \mathtt{.ok}\ c \;\wedge\;
\mathrm{Bnd56}(c) \;\wedge\; \denote{c} = \denote{a} - \denote{b},
\]
resting on the lemma $\denote{16P} = 16p \equiv 0 \pmod p$ (the
telescope, scaled by $16$), which justifies
$\denote{a + (16P - b)} = \denote{a} + 0 - \denote{b}$. Why $16$ and not
$p$ itself (i.e.\ $k = 1$): positivity of every limb of $kP - b$ requires
the \emph{smallest} limb constant, $k(2^{51}-19)$, to dominate the
\emph{largest} allowed limb of $b$, $2^{54}-1$; $k = 1$ gives
$2^{51}-19 \ll 2^{54}$ --- hopeless --- and even $k = 8$ falls $151$
short, as the worked example computed. First power of two that clears
the bar: $16$.

\solhead{10.2}
\pathway Squaring doubles the exponent; multiplying adds. Track only
exponents and check each claimed identity by substitution --- the field
elements are irrelevant to the bookkeeping, which is the insight that
makes the whole verification tractable.
\answer $a^2$: exponent $2$ ✓ (one squaring).
$a^9 = (a^2)^{2\cdot 2} \cdot a$: squaring $z_2$ twice gives exponent
$2 \cdot 2 \cdot 2 = 8$; multiplying by $a$ adds $1$: exponent $9$ ✓
(two squarings, one multiplication).
$a^{11} = a^9 \cdot a^2$: $9 + 2 = 11$ ✓ (one multiplication).
The induction that carries the rest: each ladder stage claims
$(2^j - 1)\cdot 2^k + (2^k - 1) = 2^{j+k} - 1$ for the block sizes
$(j,k)$ it uses --- verify it once in general:
$(2^j - 1)2^k + 2^k - 1 = 2^{j+k} - 2^k + 2^k - 1 = 2^{j+k} - 1$ ✓ ---
and every remaining line of the chain is an instance. Addition-chain
verification in full: substitute, cancel, done.

\solhead{10.3}
\pathway Ask what \emph{differs between the two code paths}, not between
the two mathematical functions. The math of \code{square a} and
\code{mul a a} agree; the \emph{instructions} do not.
\answer Three concrete misses. (1) \code{square} is a \emph{different
routine}: it exploits symmetry ($a_i a_j$ appears twice for $i \neq j$,
so it computes $2 a_i a_j$ once), which changes the column expressions,
the constant factors, and therefore the \emph{bounds} --- the doubled
terms eat one extra headroom bit that \code{mul}'s proof never had to
account for. (2) A bug in that symmetry trick --- the classic one is a
missing factor of $2$ on exactly one cross term --- lives \emph{only} in
\code{square}'s code path; a theorem about \code{mul} quantifies over
\code{mul}'s instructions and says nothing about it. (3) Even the
\lean{.ok} clause differs: overflow behavior depends on the actual
intermediate expressions, not on the mathematical value. The one-line
summary for design reviews: \emph{we verify code paths, not functions ---
if the optimizer wrote a second path, the verifier owes a second
theorem}. (This is also exactly why the four forks each got their own
proofs: same math, four instruction streams.)

\begin{checkpoint}
You should now be able to: state the field implementation certificate and
every quantifier in it; recite the campaign order and justify why bounds
and \code{reduce} come early; tell the memory-wall story and its
decomposition moral; and audit a stranger's field-layer claim with four
pointed questions.
\end{checkpoint}
